\documentclass[10pt,a4paper]{amsart}
\usepackage[margin=19mm]{geometry}
\usepackage{amsmath,amssymb,mathtools}
\usepackage[scr=rsfs,cal=euler]{mathalfa}
\usepackage{xfrac}
\usepackage[unicode,hidelinks,bookmarks=false]{hyperref}
\newcommand{\ie}{i.e.\ }
\newcommand{\cf}{cf.\ }
\newcommand{\resp}{respectively\ }
\newcommand{\Subsection}{\subsection}
\providecommand{\nobreakdash}{\nobreak\hbox{-}}
\setlength{\emergencystretch}{2em}
\multlinegap0pt
\usepackage{bm}
\usepackage{bbm}
\usepackage{dsfont}
\usepackage{xspace}
\usepackage{cancel}
\usepackage{mathtools}
\usepackage{cleveref}
\usepackage{amsthm}
\makeatletter
\@ifundefined{theorem}{\newtheorem{theorem}{Theorem}}{}
\@ifundefined{lemma}{\newtheorem{lemma}[theorem]{Lemma}}{}
\makeatother
\makeatletter
\newcommand{\Ca}{\mathbb{C}^n}
\def\Ha{\mathbb{H}^n}
\newcommand{\R}{\ensuremath{\mathbb{R_+}}}
\expandafter\ifx\csname similartheorem\endcsname\relax
\newenvironment{similartheorem}[1]
  {\newcommand\dualnumber{\ref*{#1}}\begin{similartheorem*}}
  {\end{similartheorem*}}
\fi
\makeatother
\title[Distances in sets of positive Kor\'anyi upper density in Hei]{Distances in sets of positive Kor\'anyi upper density in Heisenberg Group}
\author{K S Senthil Raani, Rajesh K. Singh}
\date{Source: arXiv:2507.14917v1 (CC BY 4.0)}
\begin{document}
\maketitle

\noindent\emph{Source and licence.} Distances in sets of positive Kor\'anyi upper density in Heisenberg Group. Version arXiv:2507.14917v1 (CC BY 4.0), distributed under the Creative Commons Attribution 4.0 International licence, as recorded in the arXiv licence metadata for this exact version and shown on the abstract page https://arxiv.org/abs/2507.14917v1. Full rights notice: licenses/arxiv-2507.14917-CC-BY-4.0.txt.\par\smallskip

\noindent\textit{Editorial scope.} Counted unit: the Lemma whose source label is \texttt{the LFE} in the article (source0.tex lines 646--650, source label \texttt{the LFE}), delivered together with the article's own complete original proof of it (source0.tex lines 651--716, one \texttt{proof} environment of 66 lines). No mathematics is written, repaired or re-derived: statement and proof are the article's own text line for line. Required context retained uncounted: twisted convol identity (lines 401--403); Laguerre-Asymp-Our-setup (lines 486--539); less than q by mu (lines 492--494); Surface coeff. low (lines 592--598). Every definition, display and label of those lines is retained unchanged. Omitted from this package and not redistributed: 1--400, 404--485, 495--591, 599--645, 717--1546. Nothing inside a retained excerpt is omitted or altered: every retained body is delivered exactly as the article's lines are written, so no omission receipt of any kind accompanies this package. Citations: the retained excerpts carry no citation command at all, so no bibliography entry of the article is transcribed and no reference is redistributed. Reference adaptation: none was needed and none was made; every reference inside the delivered unit resolves inside it, so no unresolved reference or citation is produced. Printed slips kept literal: no printed word, symbol, hypothesis or number of the counted statement or of its proof was altered. Layout and class: the article's own class is \texttt{amsart} with packages including amssymb, latexsym, amsmath, amsfonts, ulem, hyperref, enumitem, fancyhdr, cleveref, mathrsfs, xcolor, letterspace, soul, tikz; the wrapper is the lane's pinned \texttt{amsart} editorial wrapper at 10pt on A4, which restates the theorem-like environments the retained text needs and adds the article's own macro definitions that the retained text needs (\texttt{\textbackslash Ca}, \texttt{\textbackslash Ha}, \texttt{\textbackslash R}), with extra wrapper packages (bm, bbm, dsfont, xspace, cancel, mathtools, cleveref). The delivered numbering is the unit's own. Assets: no retained span contains a figure environment, a graphics-inclusion command, a PDF-page inclusion, a table or a TikZ picture; no image file is copied into this package, the package declares no asset dependency and no image file is redistributed with it. Licence: version arXiv:2507.14917v1 of this article is distributed under the Creative Commons Attribution 4.0 International licence, as recorded in the arXiv licence metadata for this exact version and shown on the abstract page https://arxiv.org/abs/2507.14917v1. A full rights notice is delivered at licenses/arxiv-2507.14917-CC-BY-4.0.txt. Characters: every retained excerpt is pure ASCII, so the delivered engine, XeLaTeX with its default Latin Modern text font, reports no missing character.
\begin{equation}\label{twisted convol identity}
  \sum_{|\alpha|=k} \langle \widehat{f}(\lambda)^{*} \widehat{f}(\lambda) \Phi_{\alpha}^{\lambda} , \Phi_{\alpha}^{\lambda}  \rangle = (2 \pi)^{-n} |\lambda|^n \,   \| f^{\lambda} \ast_{\lambda} \varphi_{k, \lambda}^{n-1} \|_{2}^2,
\end{equation}

\begin{lemma}\label{Laguerre-Asymp-Our-setup}
    Let $n \in \mathbb{N}$ and $k \in \mathbb{N}_{0}$. Let $0 < a < b < 1$. Set $\mu := 2 (2k +n)$. Then, there exists $k_{0} \in \mathbb{N}$, large enough, such that for $k \geq k_{0}$, the following hold: 
    \begin{enumerate}
        \item[\emph{(1)}] In the Bessel regime, that is $0 \leq x \leq b \mu$, one has
    \begin{enumerate}
 \item [\emph{(\emph{a})}] \hspace*{.7mm} If $x \,  \leq 1/ \mu$, then
    \begin{equation} \label{less than q by mu}
        \frac{\Gamma(k+1)}{\Gamma(k+n)} \,  L^{n-1}_{k}(x) \, e^{- \frac{1}{2} x} \,  x^{n-1} = C \,  x^{n-1} \left[1 + O \left(  (x \mu )^{1/2} \right) \right].
    \end{equation}  
 \item [\emph{(\emph{b})}] \hspace*{.7mm} If $ 1 / \mu \leq x \, \leq b \mu$, then
    \begin{equation}\label{1 by mu to mu}
       \begin{split}
            \frac{\Gamma(k+1)}{\Gamma(k+n)} \,  L^{n-1}_{k}(x) &  \, e^{- \frac{1}{2} x} \,  x^{n-1} = C \, x^{\frac{n}{2}-1} \mu^{-\frac{n}{2}+1}  \left( \frac{  \psi( \frac{x}{\mu} )  }{  \psi^{ \prime} ( \frac{x}{\mu} )  }  \right)^{1/2} \\
            & \times \left[  \sqrt{  \frac{2}{ \pi \mu \psi({\textstyle\frac{x}{\mu}})}      } \, \cos \left(  \mu \psi(\textstyle\frac{x}{\mu}) - \frac{\pi}{2}(n-1) - \frac{\pi}{4}      \right) + O \left(  (x \mu )^{-3/4} \right) \right]\\
           & \hspace*{5em}= x^{n-1} O \left(  (x \mu )^{-\frac{2n-1}{4}} \right).
       \end{split}
    \end{equation}      
    \end{enumerate}
     \item[\emph{(2)}] In the Airy regime, that is $a \mu \leq x \leq \frac{1}{a} \mu$,  one has
     \begin{enumerate}
 \item [\emph{(\emph{a})}] \hspace*{.7mm} If $ a \mu \leq x \leq  \mu - c \, \mu^{1/3}$, $0 < c \ll 1$, then    
 \begin{equation} \label{Bessel-Airy interface}
        \begin{split}
            \frac{\Gamma(k+1)}{\Gamma(k+n)} \,   & L^{n-1}_{k}(x) \, e^{- \frac{1}{2} x} \,  x^{n-1} \\
            & = (-1)^k \mu^{- \frac{4}{3} } \mu^{\frac{1}{12}} (\mu - x)^{-\frac{1}{4}}   \left[  \cos \left( {  \textstyle\frac{2}{3} |\xi(x)|^{3/2}  - \frac{\pi}{4} }     \right) + O \left( { \textstyle  \mu^{  \frac{1}{6}} \left( 1 - \frac{x}{\mu}  \right)^{- \frac{1}{2}} } \right)  \right],
        \end{split}
    \end{equation} 
    Here, $\xi(x) :=  - \mu^{2/3} \phi({\textstyle \frac{x}{\mu}})$. Moreover,    $ \xi(x) \sim - \mu^{2/3} \left( 1 - {\textstyle \frac{x}{\mu}}  \right)$, if $a \mu \leq x \leq \mu$ and $ \xi(x) \sim  \mu^{2/3} \left(  {\textstyle \frac{x}{\mu}}  -1 \right)$, if $x > \mu$.
 \item [\emph{(\emph{b})}] \hspace*{.7mm} If $ \mu - c \, \mu^{1/3} \leq x \leq \mu + c \, \mu^{1/3}$, $0 < c \ll 1$, then
    \begin{equation} \label{around mu}
        \begin{split}
            \frac{\Gamma(k+1)}{\Gamma(k+n)} \,   & L^{n-1}_{k}(x) \, e^{- \frac{1}{2} x} \,  x^{n-1} \\
            & = \, (-1)^{k} \mu^{{-\textstyle\frac{n}{2}} - \frac{7}{12} }  x^{ {\textstyle\frac{n}{2} - \frac{3}{4}} } \left[ \emph{Ai}( \xi(x))   +  O \left( {\textstyle \frac{1}{x} } \, \emph{Ai}( \xi(x)) \right)  \right]\\
            &= O \big(  \mu^{{-4/3 }  } \big).
        \end{split}
    \end{equation}  
    
 \item [\emph{(\emph{c})}] \hspace*{.7mm} If $ \mu + c \, \mu^{1/3} \leq x \, \leq {\textstyle \frac{3}{2}} \mu$, then
    \begin{equation}\label{Airy merging with exponential}
        \begin{split}
            \frac{\Gamma(k+1)}{\Gamma(k+n)} \,   & L^{n-1}_{k}(x) \, e^{- \frac{1}{2} x} \,  x^{n-1} \\
            & = \, (-1)^{k} \mu^{{-\textstyle\frac{4}{3}} } \,   \mu^{ {\textstyle\frac{1}{12} } }  (x- \mu)^{ - {\textstyle\frac{1}{4}  } } \, e^{ - {\frac{2}{3}  } \mu^{- \frac{1}{2}} (x- \mu)^{3/2}    } \left[   1 + O \left(   \mu^{-\frac{1}{4}} (x - \mu )^{- \frac{3}{4}}  \right)    \right].
        \end{split}
    \end{equation}  
    \end{enumerate}
    \item[\emph{(3)}] In the exponential regime, that is $x \geq \frac{3}{2} \mu$, one has
     \begin{equation}\label{expo regime}
        \begin{split}
            \frac{\Gamma(k+1)}{\Gamma(k+n)} \,   & L^{n-1}_{k}(x) \, e^{- \frac{1}{2} x} \,  x^{n-1} = O \left( \mu^{ -\frac{n}{2}-\frac{1}{2}} x^{\frac{n}{2}-1} e^{- \frac{2}{3} x}  \right).
        \end{split}
    \end{equation} 
\end{enumerate}
Here, the implied constants in the asymptotic notation may depend on $n, k_{0}$. 
\end{lemma}

\begin{lemma}\label{Surface coeff. low}
    Let $\delta \ll 1$. If $\mu |\lambda| \leq \delta \ll 1$ then one has 
    \begin{equation}\label{Low frequency surface coeff}
        R_{k}(\lambda, \sigma) =  C_{n} + O \big( (\mu |\lambda|)^{1/2} \big),
    \end{equation}
    for some $C_{n}>0$, where the implied constant in the asymptotic notation  depends on $n$.
\end{lemma}

\begin{lemma}\label{the LFE} For $1 \gg \delta \geq r^{2}$, there exists an absolute constant $C_{n}$ such that
\begin{equation}\label{low frequency contri}
    \int_{\R^{*}} \sum_{k \, : \, 2(2k +n) |\lambda | r^2 \leq \delta } R_{k}(\lambda r^{2},\sigma) \sum_{|\alpha|=k} \langle \hat{f}(\lambda)^{*} \hat{f}(\lambda) \Phi_{\alpha}^{\lambda},\Phi_{\alpha}^{\lambda} \rangle \ |\lambda|^n d\lambda \geq C_{n} |E|^2.
\end{equation}
\end{lemma}

\begin{proof}
Let us abbreviate  $\mu = 2(2k +n)$. We recall that $f= 1_{E}$ as well as the identity 
\begin{equation}\label{twisted convol identity1}
  \sum_{|\alpha|=k} \langle \widehat{f}(\lambda)^{*} \widehat{f}(\lambda) \Phi_{\alpha}^{\lambda} , \Phi_{\alpha}^{\lambda}  \rangle = (2 \pi)^{-n} |\lambda|^n \,   \| f^{\lambda} \ast_{\lambda} \varphi_{k, \lambda}^{n-1} \|_{2}^2,
\end{equation}
which is just (\ref{twisted convol identity}).

Note that,  in view of \Cref{Surface coeff. low} (applied at $\lambda r^2$ instead of $\lambda$) and the identity \eqref{twisted convol identity1}, there exists a positive constant $C_{n}$ so that
\begin{eqnarray}
\nonumber && \left|\int_{\R^{*}} \sum_{k \, : \, 2(2k +n) |\lambda | r^2 \leq \delta} \left(R_{k}(\lambda r^{2},\sigma) -C_n\right)\sum_{|\alpha|=k} \langle \hat{f}(\lambda)^{*} \hat{f}(\lambda) \Phi_{\alpha}^{\lambda},\Phi_{\alpha}^{\lambda} \rangle \ |\lambda|^n d\lambda\right|\\
&&\leq C_{n} \sqrt{\delta} \int_{\R^{*}} \sum_{k \, : \, 2(2k +n) |\lambda | r^2 \leq \delta}\| f^{\lambda} \ast_{\lambda} \varphi_{k, \lambda}^{n-1} \|_{2}^2 \ |\lambda|^{2n} d\lambda \label{eq:CompuLowThorough}
\end{eqnarray}
Hence, with an appropriate choice of $\delta<\frac14$, we see from  \eqref{twisted convol identity1} and \eqref{eq:CompuLowThorough} that the left-hand  side of \eqref{low frequency contri} is  bounded below by
    \begin{equation}\label{low frequency contri 1}
    \frac12C_n\int_{\R^{*}} \sum_{k \, : \, \mu |\lambda | r^2 \leq \delta } \,   \| f^{\lambda} \ast_{\lambda} \varphi_{k, \lambda}^{n-1} \|_{2}^2 \ |\lambda|^{2n} d\lambda.
\end{equation}
Let $\lambda>0$. We now concentrate on the factor $\| f^{\lambda} \ast_{\lambda} \varphi_{k, \lambda}^{n-1} \|_{2}^2$ in the above expression. For this we shall prove that if  $\mu \lambda \ll 1$ for $k \geq k_{0}$, with $k_{0}$ being a large natural number as in \Cref{Laguerre-Asymp-Our-setup}, and $| z| \leq \rho_{k}$ (for some $\rho_{k} \ll 1$ to be optimized later) we have
\begin{equation}\label{twisted term near 0}
    f^{\lambda} \ast_{\lambda} \varphi_{k, \lambda}^{n-1} \, (z) = C_{n}^{\prime} \, \mu^{n-1} |E| +   O (  \mu^{n-1} (\lambda \mu)^{1/2} |E|).
\end{equation}
%Then there exists an absolute non-zero constant  $C_{n}'$ such that 
%\begin{eqnarray}
%      \left|\left|f^{\lambda} \ast_{\lambda} \varphi_{k, \lambda}^{n-1} \, (z)\right|^2 - C_n^2\mu^{2n-2}|E|^2 \right|&\leq& C_{n}' \mu^{2n-2}(\lambda\mu)^{1/2}|E|^2\left((\lambda\mu)^{1/2} + 2C_n\right).\label{eq:CompuLowthorough2}
% \end{eqnarray}
Note that for $\delta \geq r^2$ the approximation \eqref{twisted term near 0} gives the lower bound for (\ref{low frequency contri 1}). Indeed, for appropriate positive constants $a_n,b_n$,  the expression \eqref{low frequency contri 1} is 


% Note that for $\delta \geq r^2$ the approximation \eqref{twisted term near 0} gives the lower bound for (\ref{low frequency contri 1}). Indeed, if $\lambda \mu \leq b$ for a sufficiently small constant $b>0$, depending on $n$ and $k_{0}$, determined by \eqref{twisted term near 0} then after restricting attention to the $k \geq k_{0}$ terms only we have (\ref{low frequency contri 1}) is
\begin{equation*}
    \begin{split}
      & \geq \frac{1}{2} C_{n} \int_{\R^{*}} \sum_{ k \geq k_{0} \, : \,  \mu \lambda \ll 1  } \,   \int_{z \in \Ca : |z | \leq \rho_{k}} |f^{\lambda} \ast_{\lambda} \varphi_{k, \lambda}^{n-1} \, (z)|^2 dz \, \lambda^{2n} d\lambda   \\
       & \geq \frac18 C_{n}{C_n^{\prime}}^2 |E|^2 \,   \int_{0 <\lambda \,  \leq  b_n} \sum_{k_{0} \leq  k \leq a_{n} \lambda^{-1}   } \,   \rho_{k}^{2n} \mu^{2n-2} \ \lambda^{2n} d\lambda.
    \end{split}
\end{equation*}
Now $\rho_{k}$ can be chosen to be $\mu^{-M}$, for some $M>0$, and the last expression then becomes
\begin{equation*}
    \begin{split}       
       &  \frac18 C_{n}{C_n^{\prime}}^2 |E|^2 \,   \int_{0 <\lambda \,  \leq  b_n} \sum_{k_{0} \leq  k \leq a_{n} \lambda^{-1}   } \,  \mu^{-2n M + 2n-2} \ \lambda^{2n} d\lambda,
\end{split}
\end{equation*}
which is $\simeq_{n, k_{0}} |E|^2$. This gives the asserted lower bound for \eqref{low frequency contri 1}, and hence \eqref{low frequency contri} is proved. 

It remains to establish \eqref{twisted term near 0}. Here we use asymptotic of Laguerre polynomials in the Bessel regime and write  
    \begin{equation*}\label{3 twisted term near 0}
    \begin{split}
        f^{\lambda} \ast_{\lambda} \varphi_{k, \lambda}^{n-1} \, (z) &  = \int_{\mathbb H^n} 1_{E}(z-w,-t) \, L_k^{n-1}\big( {\textstyle \frac{1}{2}} \lambda|w|^2 \big) \, e^{- \frac{1}{4}\lambda |w|^2}   e^{i \lambda \left(-t+\frac12 \Im(z\cdot \overline{w}) \right)}dwdt \\
        & = \int_{(z,0) - E \,  \subseteq \mathbb H^n}  L_k^{n-1}\big( {\textstyle \frac{1}{2}} \lambda|w|^2 \big) \, e^{- \frac{1}{4}\lambda |w|^2}   e^{i \lambda \left(-t+\frac12 \Im(z\cdot \overline{w}) \right)}dwdt
    \end{split}
\end{equation*}
where,  we have used that  $1_{E}(z-w, -t)=1$ if and only if $(w,t) \in (z,0) - E$. Since $E \subseteq B(0,1) \subset \Ha$, the Kor\'anyi norm $|(w-z,t)|_{K} <1$, whence  $\frac{\lambda}{4} |w |^2 \leq \lambda$ and $| \lambda \left(-t+\frac12 \Im(z\cdot \overline{w}) \right) | \leq 2\lambda$. Since $\lambda \ll  1$, so $ \left| \exp{ i \lambda \left(-t+\frac12 \Im(z\cdot \overline{w}) \right) } - 1 \right| \leq 4 \lambda$ which together with the estimate $\frac{\Gamma(k+1)}{\Gamma(n)\Gamma(k+n)} |L_{k}^{n-1}(x) e^{-x/2}| \leq 1$,  and the Stirling approximation, gives that the last integral equals
\begin{equation}\label{4 twisted term near 0}
    \begin{split}
         \int_{(z,0) - E \,  \subseteq \mathbb H^n}  L_k^{n-1}\big( {\textstyle \frac{1}{2}} \lambda|w|^2 \big) \, e^{- \frac{1}{4}\lambda |w|^2}  dwdt \ + \ O \left( \lambda \mu^{n-1} |E| \right).
    \end{split}
\end{equation}
Since, ${\textstyle\frac{1}{2}} \lambda |w|^2 \leq 2 \lambda \ll 1/ \mu$  in the integral in (\ref{4 twisted term near 0}), using Laguerre polynomial asymptotic (\ref{less than q by mu}) in conjugation with the Stirling approximation the last expression (\ref{4 twisted term near 0})  becomes
\begin{equation*}\label{5 twisted term near 0}
    \begin{split}
        & C  \int_{(z,0) - E \,  \subseteq \mathbb H^n} \frac{\Gamma(k+n)}{\Gamma(k+1)} \left[ 1 + O  \left( \left( {\textstyle\frac{1}{2}} \lambda |w|^2 \, \mu \right)^{1/2} \right) \right]  dwdt \ + \ O \left( \lambda \mu^{n-1} |E| \right)\\
        & \simeq \mu^{n-1} |E| +  O (  \mu^{n-1} (\lambda \mu)^{1/2} |E|),
    \end{split}
\end{equation*}
which is what we were aiming for in (\ref{twisted term near 0}). Thus, we have completed the proof of \Cref{the LFE}.
 % \blue{On the other hand, for $0 \leq k \leq k_{0}$, we argue as in \Cref{Surface coeff. low}. Thus, for  Lguerre polynomials, $L_{k}^{n-1}$, with their degrees bounded by $k_{0}$, one has $L_{k}^{n-1}(x) = \frac{\Gamma(k+n)}{\Gamma(k+1) \Gamma(n)} + O_{k_{0}}(x)$ valid for small $x$ unformly for all $0 \leq k \leq k_{0}$. Substituting this into (\ref{4 twisted term near 0}), one has that for $\lambda$ sufficiently small the left side of (\ref{4 twisted term near 0}) becomes $\simeq_{k_{0}} |E| + O_{k_{0}}(\lambda|E|)$.} \red{(Writing of remaining portion of bounded $k$'s can only start when the writing of large $k$ is done!!!)}
 % Thus, we have completed the proof of \Cref{the LFE}.
\end{proof}

\end{document}
